% !TEX root = ../main.tex
\chapter{Introduction}
\begin{justify}
This chapter introduces CrickEye — a low-cost, computer-vision-based cricket coaching system for batting net sessions — and summarises motivation, problem, objectives, scope, and report structure.

\section{Overview}

\par Cricket in India spans structured pathways under the BCCI and state associations \cite{BCCIportal}. Large ICC events underline economic scale: the ICC cites a Nielsen assessment of the 2023 Men's World Cup in India at roughly USD~1.39 billion total economic impact \cite{ICC2023WorldCupIndia}.

\par Surveys of computer vision in sports note that production-grade analytics historically assumed multi-camera rigs and controlled conditions \cite{Thomas2017}, leaving smaller academies underserved --- the gap CrickEye targets.

\par Coaching infrastructure is uneven. Metro academies use video suites, Hawk\-Eye\-class tracking, high-speed cameras, wearables, and biomechanics support. Tier-2 and Tier-3 centres often rely on one coach observing many batters with little or no instrumentation.

\par The resulting gap is less about talent than about access to repeatable, quantitative feedback.

\section{Problem Statement}

\par The core problem addressed by this project is: \textit{How can a cricket academy with limited resources provide systematic, quantitative, and repeatable biomechanical feedback to its batters using nothing more than a standard smartphone camera?}

\par Commercial options are either costly instrumented stacks (radar, Hawk-Eye-class systems, wearables) or generic video tools that depend on manual tagging. Neither fits a typical academy that wants automated feedback from one front-on commodity camera. CrickEye targets that niche.

\section{Objectives}

\par The objectives of this project are:

\begin{enumerate}
    \item To design an end-to-end pipeline on a single front-on net video: detect swings, classify shots, and score biomechanics per delivery.
    \item To implement pose-based swing detection with YOLOv8 on continuous net footage.
    \item To train Swin3D-T on five shot classes (cover, flick, pull, straight, sweep) from public YouTube clips.
    \item To extract front-on-safe signals: head, feet, stance symmetry, elbow category, swing-path proxy, plus diagnostic swing intensity and bat speed.
    \item To map metrics to a 0--10 composite (raw weights 28/20/12/12/22/6) and surface repeated-pattern coaching alerts.
    \item To add optional ball tracking for speed, bounce, and length zones.
    \item To ship a web dashboard (WebSockets, per-shot cards, wagon wheel, pitch map, player/coach roles).
    \item To validate on held-out YouTube nets and commodity hardware.
\end{enumerate}

\section{Scope of the Project}

\par CrickEye is designed for the specific use case of analysing batting net sessions filmed from a \textbf{front-on} (batter-facing) camera angle. The system is not intended for broadcast-quality multi-camera analysis or for bowling/fielding analysis. The key boundaries of the current scope are:

\begin{itemize}
    \item \textbf{Camera:} Single, stationary, front-on camera (e.g. a smartphone on a tripod).
    \item \textbf{Subject:} One primary batter per session; the system selects the largest detected person in each frame.
    \item \textbf{Shot types:} Five classes (cover, flick, pull, straight, sweep).
    \item \textbf{Metrics:} Composite uses only front-on-safe 2D cues, plus execution when ball length exists. Depth-heavy quantities (weight transfer, spine angle, knee flex) may log but do not enter the headline score.
    \item \textbf{Diagnostics:} Swing intensity and bat speed are trend and sanity channels, not composite weights.
    \item \textbf{Ball tracking:} Optional module using YOLO-based ball detection; speed and length estimates are single-camera approximations, not radar-grade measurements.
    \item \textbf{Deployment:} Local web application; cloud integration via Supabase for authentication and session persistence.
\end{itemize}

\section{Significance of the Study}

\par The work contributes a practical, low-hardware path to AI-assisted batting coaching in India. Contributions:

\begin{enumerate}
    \item An integrated pipeline: pose estimation plus a 3D transformer for cricket shot recognition.
    \item Front-on-safe biomechanics with quality cues, shot-type-aware rules, optional execution from ball length, and diagnostic swing-speed channels.
    \item A coaching-oriented dashboard (alerts, drills, session trends) that speaks in coach-facing language.
    \item A YouTube-sourced training recipe showing competitive five-way classification without proprietary data.
\end{enumerate}

\section{Organisation of the Report}

\par The remainder of this report is organised as follows:

\begin{itemize}
    \item \textbf{Chapter 2: Literature Review} --- Reviews existing work in computer vision for sports, pose estimation, video classification, and cricket-specific analytics.
    \item \textbf{Chapter 3: Methodology} --- Describes the system architecture, pipeline design, dataset collection, model training, biomechanical metric extraction, scoring logic, and frontend/backend implementation.
    \item \textbf{Chapter 4: Results} --- Outputs, metrics, classifier evaluation, dashboard illustrations, limitations, and (per MIET format) \textbf{Conclusion} and \textbf{Future Directions} in the same chapter.
\end{itemize}

\end{justify}

